% 第9章第5节：一般方阵特征子空间的扰动（不变子空间的分块形式，特征方向的偏转界，非正规性放大偏转）。

\section{一般方阵特征子空间的扰动}
\label{sec:ls-nonnormal-subspace-sensitivity}

对称矩阵的特征子空间两两正交，第~\ref{sec:ls-subspace-sensitivity}节的偏转只由扰动与谱隙之比控制。一般方阵的特征向量彼此不正交，上一节已经看到，这会放大特征值的误差。本节说明特征方向也是如此：所选方向的偏转除了取决于特征值之间的间隔，还取决于其余部分的非正规程度。

为了只用实数，本节考虑实特征值与实特征向量；复特征值的情形可以照样在复数域中讨论。

\subsection{不变子空间的分块形式}
一般方阵的特征向量不一定能组成基，所以比较的对象取为不变子空间。设$\bm A\in\mathbb R^{n\times n}$，$\mathcal S$是$\bm A$的$k$维不变子空间，即$\bm A\mathcal S\subseteq\mathcal S$，例如一个实特征值的特征空间。取$\bm Q_1\in\mathbb R^{n\times k}$的列为$\mathcal S$的标准正交基，再补成正交矩阵$\bm Q=[\bm Q_1,\bm Q_2]$。由于$\bm A\bm Q_1$的列仍在$\mathcal S$中，与$\bm Q_2$的列正交，所以$\bm Q_2^\top\bm A\bm Q_1=\bm0$，于是
\[
 \bm Q^\top\bm A\bm Q=\begin{bmatrix}\bm A_{11}&\bm A_{12}\\\bm0&\bm A_{22}\end{bmatrix},\qquad
 \bm A_{11}=\bm Q_1^\top\bm A\bm Q_1,\quad\bm A_{22}=\bm Q_2^\top\bm A\bm Q_2 .
\]
分块上三角矩阵的特征多项式是两个对角块特征多项式之积，所以$\bm A_{11}$的特征值是$\bm A$限制在$\mathcal S$上的特征值，$\bm A_{22}$的特征值是其余的特征值。

$\bm A$对称时，$\bm A_{12}=(\bm Q_2^\top\bm A\bm Q_1)^\top=\bm0$，矩阵化为分块对角形，$\bm A_{22}$也对称。一般方阵的$\bm A_{12}$不必为零，$\bm A_{22}$也不必正规。下面会看到，偏转界中出现的正是$\bm A_{22}$。

一个向量$\bm w\neq\bm0$与$\mathcal S$的夹角$\theta$满足$\sin\theta=\lVert(\bm I-\bm Q_1\bm Q_1^\top)\bm w\rVert_2/\lVert\bm w\rVert_2$，也就是$\bm w$投影不到$\mathcal S$中的部分所占的比例。由$\bm I-\bm Q_1\bm Q_1^\top=\bm Q_2\bm Q_2^\top$与$\bm Q_2$列标准正交，
\[
 \sin\theta=\frac{\lVert\bm Q_2^\top\bm w\rVert_2}{\lVert\bm w\rVert_2}.
\]
$k=1$时，这就是第~\ref{sec:ls-subspace-sensitivity}节两条直线之间的主夹角。

\subsection{特征方向的偏转界}
设$\mu$是$\bm A+\bm E$的实特征值，$\bm w$是对应的实特征向量。我们关心$\bm w$离$\mathcal S$有多远。若$\mu$不是$\bm A_{22}$的特征值，答案由$\bm A_{22}-\mu\bm I$的逆控制。
\begin{theorem}[特征方向偏离不变子空间的界]
\label{thm:ls-invariant-subspace-deflection}
在上述记号下，设$(\bm A+\bm E)\bm w=\mu\bm w$，$\bm w\neq\bm0$，且$\bm A_{22}-\mu\bm I$可逆。则$\bm w$与$\mathcal S$的夹角$\theta$满足
\[
 \sin\theta
 =\frac{\lVert(\bm A_{22}-\mu\bm I)^{-1}\bm Q_2^\top\bm E\bm w\rVert_2}{\lVert\bm w\rVert_2}
 \leq\lVert(\bm A_{22}-\mu\bm I)^{-1}\rVert_2\,\lVert\bm E\rVert_2 .
\]
\end{theorem}
\begin{proof}
把特征方程写成$(\bm A-\mu\bm I)\bm w=-\bm E\bm w$，左乘$\bm Q_2^\top$。由$\bm Q_2^\top\bm A=\bm Q_2^\top\bm A(\bm Q_1\bm Q_1^\top+\bm Q_2\bm Q_2^\top)=\bm A_{22}\bm Q_2^\top$（这里用到$\bm Q_2^\top\bm A\bm Q_1=\bm0$），得
\[
 (\bm A_{22}-\mu\bm I)\bm Q_2^\top\bm w=-\bm Q_2^\top\bm E\bm w ,\qquad
 \bm Q_2^\top\bm w=-(\bm A_{22}-\mu\bm I)^{-1}\bm Q_2^\top\bm E\bm w .
\]
两边取范数并除以$\lVert\bm w\rVert_2$，由$\sin\theta=\lVert\bm Q_2^\top\bm w\rVert_2/\lVert\bm w\rVert_2$得等式。再由$\lVert\bm Q_2^\top\bm E\bm w\rVert_2\leq\lVert\bm E\bm w\rVert_2\leq\lVert\bm E\rVert_2\lVert\bm w\rVert_2$得不等式。
\end{proof}

这个界只要求$\mu$不是其余特征值，不要求$\bm A$可对角化，也不要求$\mathcal S$中的特征值与其余特征值保持某种次序。剩下的问题是$\lVert(\bm A_{22}-\mu\bm I)^{-1}\rVert_2$有多大。

对称情形下，$\bm A_{22}$对称，它的特征值$\lambda_j$就是$\bm A$在$\mathcal S$之外的特征值。对称矩阵$\bm A_{22}-\mu\bm I$的逆的谱范数是$1/\min_j|\lambda_j-\mu|$，所以
\[
 \sin\theta\leq\frac{\lVert\bm E\rVert_2}{\min_j|\lambda_j-\mu|},
\]
分母是扰动后的特征值$\mu$与组外特征值之间的距离，这正是第~\ref{sec:ls-subspace-sensitivity}节交叉间隔$\gamma$的作用。

一般情形下，若$\bm A_{22}$可对角化，$\bm A_{22}=\bm W\bm D\bm W^{-1}$，则$(\bm A_{22}-\mu\bm I)^{-1}=\bm W(\bm D-\mu\bm I)^{-1}\bm W^{-1}$，所以
\[
 \lVert(\bm A_{22}-\mu\bm I)^{-1}\rVert_2\leq\frac{\kappa_2(\bm W)}{\min_j|\lambda_j-\mu|}.
\]
特征值间隔仍在分母上，但前面多了特征基的条件数$\kappa_2(\bm W)$，与Bauer--Fike界中的因子同一来源（$\bm W$可能需要取复矩阵，这时按上一节的复欧氏范数理解）。$\bm A_{22}$不可对角化或接近不可对角化时，$\lVert(\bm A_{22}-\mu\bm I)^{-1}\rVert_2$可以远大于$1/\min_j|\lambda_j-\mu|$。

\subsection{非正规性放大偏转}
下面的例子中，特征值间隔始终为$1$，所选特征值在扰动下完全不动，特征方向却按非正规程度成比例地偏转。
\begin{example}[特征值不动而特征方向偏转]
\label{ex:ls-nonnormal-eigenvector-deflection}
取$c\geq0$与
\[
 \bm A=\begin{bmatrix}0&0&0\\0&1&c\\0&0&1\end{bmatrix},\qquad
 \bm E=t\,\bm e_3\bm e_1^\top=\begin{bmatrix}0&0&0\\0&0&0\\t&0&0\end{bmatrix},\qquad t>0 .
\]
$\bm A$的特征值为$0,1,1$，特征值$0$的特征方向是$\bm e_1$。取$\mathcal S=\operatorname{span}\{\bm e_1\}$，$\bm Q$取单位矩阵，则
\[
 \bm A_{22}=\begin{bmatrix}1&c\\0&1\end{bmatrix},
\]
特征值$0$与组外特征值的距离为$1$，与$c$无关；$c>0$时$\bm A_{22}$不可对角化。

扰动后的矩阵第一行仍为零，特征多项式为$z(z-1)^2$，特征值$0$保持不动。解$(\bm A+\bm E)\bm w=\bm0$：第三行给出$tw_1+w_3=0$，第二行给出$w_2+cw_3=0$，取$w_1=1$得
\[
 \bm w=\begin{bmatrix}1\\ct\\-t\end{bmatrix},\qquad
 \sin\theta=\frac{t\sqrt{1+c^2}}{\sqrt{1+t^2(1+c^2)}} .
\]
$t$很小时$\sin\theta\approx t\sqrt{1+c^2}$。$\lVert\bm E\rVert_2=t$，所以特征方向偏转的量级是扰动的$\sqrt{1+c^2}$倍；$c=100$时，大小为$10^{-6}$的扰动使特征方向转过约$10^{-4}$。

与定理~\ref{thm:ls-invariant-subspace-deflection}比较：$\mu=0$，$\bm A_{22}^{-1}=\left[\begin{smallmatrix}1&-c\\0&1\end{smallmatrix}\right]$，$\bm Q_2^\top\bm E\bm w=(0,t)^\top$，所以$\bm A_{22}^{-1}\bm Q_2^\top\bm E\bm w=(-ct,t)^\top$，长度为$t\sqrt{1+c^2}$，除以$\lVert\bm w\rVert_2$恰好给出上面的$\sin\theta$，与定理中的等式一致。放大倍数来自$\lVert\bm A_{22}^{-1}\rVert_2$，它随$c$增长，而特征值间隔始终为$1$。
\end{example}
这个例子中特征值稳定、间隔不小，但所选特征方向仍然敏感。所以对一般方阵，单看特征值及其间隔不能判断特征方向的稳定性，还要检查其余部分的非正规程度，即$\lVert(\bm A_{22}-\mu\bm I)^{-1}\rVert_2$与$1/\min_j|\lambda_j-\mu|$相差多少。对称矩阵没有这个问题，奇异值分解也是如此：它把问题放进对称增广矩阵中处理，下面两节分别讨论奇异值与奇异子空间。
